\section{From square tori to the full plane}\label{red:section}\label{sec:reduction}

This section proves the reduction needed for the full admissible class.
A uniform lower bound on finite square tori gives the same lower bound
for every field in $\mathcal A_1$.  The argument uses the periodic
approximation theorem of Sandier and Serfaty, after checking the
normalization and the fixed-width cutoffs in the definition of $W$.

\subsection{Local regularity and the puncture limit}

The local integrability convention causes no restriction on the
competitors: under the distributional equations, membership in
$L^1_{\mathrm{loc}}$ already implies membership in every
$L^q_{\mathrm{loc}}$, $1\le q<2$.

\begin{lemma}\label{red:local}
Let $\Lambda\subset\R^2$ be simple and locally finite, and let
$E\in\mathcal D'(\R^2;\R^2)$ be a distributional vector field satisfying
\[
 \operatorname{div}E=2\pi\Bigl(\sum_{p\in\Lambda}\delta_p-1\Bigr),
 \qquad \operatorname{curl}E=0.
\]
The field is smooth off $\Lambda$.  Near each $p\in\Lambda$ it has
the form
\begin{equation}\label{red:local-expansion}
 E(x)=\frac{x-p}{|x-p|^2}+F_p(x),\qquad F_p\text{ smooth}.
\end{equation}
Consequently $E$ has a representative in $L^q_{\mathrm{loc}}$ for
every $1\le q<2$, and
$W(E,\xi)$ exists as a finite real number for every smooth compactly
supported $\xi$.
\end{lemma}

\begin{proof}
Choose a disk about $p$ containing no other charge and subtract
\[
 E_p^0(x)=\nabla\log|x-p|-\pi(x-p).
\]
Since $\Delta\log|x-p|=2\pi\delta_p$, the difference
$E-E_p^0$ has zero divergence and curl as a distribution on the disk.
The identities
\[
 \Delta V_1=\partial_1\operatorname{div}V-
                 \partial_2\operatorname{curl}V,
 \qquad
 \Delta V_2=\partial_2\operatorname{div}V+
                 \partial_1\operatorname{curl}V
\]
and Weyl's lemma make both components of this difference smooth.
On a disk without charges, use $E^0(x)=-\pi x$ instead.
This proves \eqref{red:local-expansion} and smoothness off the charges.
The singular function $|x-p|^{-1}$ is locally in $L^q$ exactly when
$q<2$.

Only finitely many charges meet a fixed compact neighborhood of
$\operatorname{supp}\xi$.  Choose disjoint small disks about them.
In each disk,
\[
 \frac12|E(x)|^2-\frac1{2|x-p|^2}\in L^1.
\]
Moreover $\xi(x)-\xi(p)=O(|x-p|)$, so
$(\xi(x)-\xi(p))/|x-p|^2$ is integrable.  The remaining singular
integral is $\pi\xi(p)\log(\rho/\eta)$, whose divergence cancels
the prescribed $\pi\xi(p)\log\eta$.  The integral outside these
disks is finite.  Thus the puncture limit exists and is finite.
\end{proof}

In particular the only possible infinite values of $W(E)$ arise in
the large-square limsup.  No puncture limit is interchanged with that
limsup below.

\subsection{Approximation and local energy estimates}

We distinguish the current normalization from the electric-field
normalization of $\mathcal A_1$.  For $b>0$, write $\mathcal C_b$ for
currents $j\in L^1_{\mathrm{loc}}(\R^2;\R^2)$
with
\[
 \operatorname{div}j=0,\qquad
 \operatorname{curl}j=\rho-b,\qquad
 \rho=2\pi\sum_{v\in\Lambda}\delta_v,
 \qquad \rho(B_s)=O(s^2),
\]
where $\Lambda$ is simple and locally finite.  Denote their local
puncture functional by $\mathsf W(j,\xi)$; its formula is the same
as that of $W(E,\xi)$.  Put $Q_s=(-s,s)^2$.

For such a current, the inverse component rotation $F=(j_2,-j_1)$
satisfies $\operatorname{div}F=\rho-b$ and $\operatorname{curl}F=0$.
In the proof of Lemma~\ref{red:local}, use the local model
$\nabla\log|x-p|-(b/2)(x-p)$ near a charge and $-(b/2)x$ off the
charges.  This gives the radial singular expansion for $F$ and its
rotated expansion for $j$.  Orthogonality preserves the squared norm,
so the same finite local puncture conclusion holds for $j$ and every
smooth compactly supported cutoff.  This local argument uses simplicity
and local finiteness, but not the charge-growth bound.

For balls and for $Q_s$, choose cutoffs $0\le\zeta_{U_s}\le1$,
compactly supported in $U_s$, equal to one at distance at least one
from its boundary, and with uniformly bounded gradients.  Write
\[
 \mathsf W_U(j)=\limsup_{s\to\infty}
       \frac{\mathsf W(j,\zeta_{U_s})}{|U_s|}.
\]
Theorem~1(1),(3),(4) of \cite{SandierSerfaty2012} makes the cutoff
choice immaterial within each fixed shape family, gives a common finite
minimum
\begin{equation}\label{red:M}
 M=\min_{\mathcal C_1}\mathsf W_{\mathrm{ball}}
   =\min_{\mathcal C_1}\mathsf W_Q\in\R,
\end{equation}
and provides square-periodic $j_k\in\mathcal C_1$ with
$\mathsf W_{\mathrm{ball}}(j_k)\to M$.  Their square periods have
side lengths $\sqrt{2\pi n_k}$ for positive integers $n_k$.
The common minimum does not assert equality of the ball and square
values of an arbitrary fixed field.
Their conditions (1.4)--(1.5) require increasing bounded open sets
with $\bigcap_{s>0}\overline U_s=\{0\}$, left-continuous volumes,
$|U_s-U_s|\le C|U_s|$, and, for every fixed $a\in\R^2$,
\[
 \frac{|(a+U_s)\mathbin\triangle U_s|}{|U_s|}\longrightarrow0.
\]
They further require constants $C$ and $\theta<2$ such that
\[
 U_s+B_1\subset U_{s+C},\qquad
 U_{s+1}\subset U_s+B_C,\qquad
 |U_{s+1}\setminus U_s|=O(s^\theta).
\]
Balls and squares satisfy these conditions with $\theta=1$ and
$|U_s-U_s|=4|U_s|$.

We also use the following nonsharp consequences of Lemma~4.7 and
Proposition~4.9 of \cite{SandierSerfaty2012}.
The latter is an application of the mass-displacement method developed
for Ginzburg--Landau energies in~\cite{SandierSerfatyMassDisplacement2011}.
In the proof of their Lemma~4.7, Sandier and Serfaty credit the
dyadic-shell argument for the local $L^r$ estimate to
Struwe~\cite{Struwe1994}; the precise weighted estimate used here is
theirs.
Let $U$ be bounded and open, $\widehat U=U+B_1$, and suppose
$\operatorname{div}j=0$ and $\operatorname{curl}j=\rho-1$
in $\widehat U$.  The number $n$ of simple charges there is finite,
and $j$ is locally square-integrable off those charges.
For $1<r<2$, $0\le\xi\le1$ compactly supported in $U$, and
$\|\nabla\xi\|_\infty\le C_0$, Lemma~4.7 implies
\begin{equation}\label{red:local-Lr}
 \int_U\xi^{r/2}|j|^r
 \le C_{r,C_0}(|U|+1)^{1-r/2}
 \bigl(\mathsf W(j,\xi)_++n\log(2+n)+1\bigr)^{r/2},
\end{equation}
where $a_+=\max\{a,0\}$.
For $n\ge1$, this follows by replacing the energy by its positive
part, enlarging the logarithmic count, and absorbing the fixed cutoff
bounds into the constant.  The statement of Lemma~4.7 and the final
display of its proof differ by a factor two on the energy term; this
weakening covers either form.  If $n=0$, there is no puncture term and
H\"older's inequality gives directly
\[
 \int_U\xi^{r/2}|j|^r
 \le |U|^{1-r/2}\left(\int_U\xi|j|^2\right)^{r/2}
 =|U|^{1-r/2}\bigl(2\mathsf W(j,\xi)\bigr)^{r/2},
\]
which also implies \eqref{red:local-Lr}.
Proposition~4.9 provides one measure $g_U$, supported in $\widehat U$
and satisfying $g_U\ge-C\,\dd x$, such that
for every smooth compactly supported $\xi$ in $U$,
\begin{equation}\label{red:mass-displacement}
 \left|\mathsf W(j,\xi)-\int\xi\,\dd g_U\right|
 \le C n_\xi\bigl(\log(2+n_\xi)+1\bigr)
                  \|\nabla\xi\|_\infty,
\end{equation}
where
\[
 n_\xi=\#\{v\in\Lambda:
       B(v,C)\cap\operatorname{supp}\nabla\xi\ne\varnothing\}.
\]
The constants in \eqref{red:mass-displacement} are universal.
In particular the same $g_U$ works for two different cutoffs.
The local smoothness needed in these results follows as in
Lemma~\ref{red:local}, after rotating the components.

\subsection{Scaling and comparison of fixed-width cutoffs}

\begin{lemma}\label{red:plane-bound}
Every $E\in\mathcal A_1$ satisfies
\begin{equation}\label{red:plane-lower}
 W(E)\ge 2\pi\left(M-\frac14\log(2\pi)\right).
\end{equation}
In particular $W(E)\ne-\infty$.
\end{lemma}

\begin{proof}
The case $W(E)=+\infty$ is immediate.  Assume that its limsup is
bounded above, allowing provisionally $W(E)=-\infty$.
Set $m=2\pi$ and rotate the vector components by
\[
 j=(-E_2,E_1),\qquad
 j'(y)=m^{-1/2}j(y/\sqrt m),\qquad
 \xi_t(y)=\chi_{t/\sqrt m}(y/\sqrt m).
\]
Thus $j\in\mathcal C_m$ and $j'\in\mathcal C_1$;
the charges of $j'$ are at $\sqrt m\,\Lambda$.
The cutoff $\xi_t$ is supported in $\overline Q_t$, equals one on
$\overline Q_{t-\sqrt m}$, and has a uniformly bounded gradient.
Changing variables in the punctured integral changes its puncture
radius from $\eta$ to $\sqrt m\,\eta$.  Consequently, at each fixed
$t=\sqrt m R$,
\begin{equation}\label{red:exact-scaling}
 W(E,\chi_R)=\mathsf W(j',\xi_t)
       -\frac\pi2\log m\sum_{v\in\sqrt m\Lambda}\xi_t(v).
\end{equation}
There is no area factor in this local identity.

The charge-growth hypothesis and the upper bound on the limsup give
$\mathsf W(j',\xi_t)\le Ct^2$ for all sufficiently large $t$.
For any fixed $A>0$, apply \eqref{red:local-Lr} to $\xi_{t+D}$
with a fixed $D>A+\sqrt m+2$, in $U=Q_{t+D+1}$.
There are $O(t^2)$ charges in $\widehat U$, and this cutoff is one
on $Q_{t+A}$, so
\begin{equation}\label{red:bulk-Lr}
 \int_{Q_{t+A}}|j'|^r\le C_{r,A}t^2(\log t)^{r/2},
 \qquad 1<r<2.
\end{equation}

Here are the charge estimates that make the boundary error negligible.
Fix $A>0$.  Coarea for $y\mapsto\|y\|_\infty$ and
\eqref{red:bulk-Lr} with $A+2$ in place of $A$ give radii
\[
 s_+\in(t+A,t+A+1),\qquad
 s_-\in(t-A-1,t-A)
\]
such that both boundary integrals of $|j'|^r$ are at most
$Ct^2(\log t)^{r/2}$.  They may be chosen with no charge on either
boundary.  Stokes' theorem, first on the punctured square and then
letting the punctures shrink, yields
\[
 \rho'(Q_{s_\pm})-|Q_{s_\pm}|
       =\int_{\partial Q_{s_\pm}}j'\cdot\tau\,\dd\mathcal H^1,
 \qquad \rho'=2\pi\sum_{v\in\sqrt m\Lambda}\delta_v.
\]
H\"older's inequality bounds the absolute value by
$Ct^{1+1/r}(\log t)^{1/2}$.  Monotonicity of $\rho'$ therefore gives
\begin{align}
 \rho'(Q_t)&=|Q_t|+o(t^2),\label{red:density}\\
 \rho'(Q_{t+A}\setminus Q_{t-A})
 &\le C_A\bigl(t+t^{1+1/r}(\log t)^{1/2}\bigr)
   =o(t^2/\log t).\label{red:strip}
\end{align}
The last equality uses $r>1$.  In particular
\begin{equation}\label{red:weighted-density}
 \sum_v\xi_t(v)=\frac{|Q_t|}{2\pi}+o(t^2).
\end{equation}

Choose a fixed $L>\sqrt m+2$ and a standard unit-width cutoff
$\zeta_{Q_{t-L}}$.  Both this cutoff and $\xi_t$ are compactly
supported in $U=Q_{t+1}$.  Their gradient supports, together with
any fixed-radius neighborhoods, lie in a strip of the form appearing
in \eqref{red:strip}.  Hence each error in
\eqref{red:mass-displacement} is $o(t^2)$.
Since $0\le\zeta_{Q_{t-L}}\le\xi_t\le1$, and their difference is
supported in $\overline Q_t\setminus Q_{t-L-1}$,
\[
 \int(\xi_t-\zeta_{Q_{t-L}})\,\dd g_U
 \ge-C\int(\xi_t-\zeta_{Q_{t-L}})\,\dd x\ge-O(t).
\]
Using this one measure for both cutoffs proves
\begin{equation}\label{red:cutoff-comparison}
 \mathsf W(j',\xi_t)
 \ge\mathsf W(j',\zeta_{Q_{t-L}})-o(t^2).
\end{equation}
Every boundary width used here is fixed as $t\to\infty$.

Divide \eqref{red:cutoff-comparison} by $|Q_t|$ and take limsups.
The ratio $|Q_{t-L}|/|Q_t|$ tends to one, so
\[
 \limsup_{t\to\infty}\frac{\mathsf W(j',\xi_t)}{|Q_t|}
 \ge\mathsf W_Q(j')\ge M.
\]
These inequalities hold in the extended real sense.  From
\eqref{red:exact-scaling}, \eqref{red:weighted-density}, and
$|Q_t|=m|K_R|$, we obtain
\[
 W(E)=m\left(\limsup_{t\to\infty}
                \frac{\mathsf W(j',\xi_t)}{|Q_t|}
                         -\frac14\log m\right).
\]
This proves \eqref{red:plane-lower} and rules out the provisional
case $W(E)=-\infty$.
\end{proof}

\subsection{Periodic averaging and the harmonic component}

For a periodic field $F$ with one unit logarithmic singularity at
each of $N$ charges in a period cell $P$, define its cell energy by
\[
 \mathcal E_P(F)=\lim_{\eta\downarrow0}
 \left[\frac12\int_{P\setminus\bigcup_v B(v,\eta)}|F|^2\,\dd x
                         +\pi N\log\eta\right].
\]
The integral is interpreted on the flat torus, so disks meeting a
chosen cell boundary are identified periodically.

\begin{lemma}\label{red:periodic-average}
Let $F$ be a periodic admissible electric field or current, with
simple charges and period cell $P$ of area $A$.  For either square
or ball cutoffs with fixed boundary width and uniformly bounded
gradients, its energy per unit area has a limit, equal to
$\mathcal E_P(F)/A$.
\end{lemma}

\begin{proof}
There are finitely many charges on the fixed torus.  Choose
$\rho>0$ smaller than its injectivity radius, with the disks about
distinct torus charges disjoint, and write on its periodic lift
\[
 F_\rho(x)=\frac12|F(x)|^2-
       \sum_v\frac{\mathbf1_{B(v,\rho)}(x)}{2|x-v|^2}.
\]
By the local expansion in Lemma~\ref{red:local}, with rotated
components for currents, $F_\rho$ is periodic and belongs to
$L^1(P)$.  For a cutoff $\xi$ the puncture calculation gives
\begin{align*}
 \mathsf W(F,\xi)
 &=\int\xi F_\rho\,\dd x+\pi\log\rho\sum_v\xi(v)+\mathcal R_\xi,\\
 \mathcal R_\xi
 &=\sum_v\int_{B(v,\rho)}
             \frac{\xi(x)-\xi(v)}{2|x-v|^2}\,\dd x.
\end{align*}
Here $\mathsf W$ denotes the same local formula for either type of
field.  A summand in $\mathcal R_\xi$ vanishes unless
$B(v,\rho)$ meets $\operatorname{supp}\nabla\xi$, and each has
absolute value at most $\pi\rho\|\nabla\xi\|_\infty$.
For a ball or square of size $R$ there are $O(R)$ such charges,
with constants allowed to depend on the fixed periodic field.
Thus $\mathcal R_\xi=O(R)$.

The number of period cells meeting any fixed-width boundary strip
is also $O(R)$.  Integrating the periodic $L^1$ function
$F_\rho$ over the remaining complete cells, and bounding the
boundary cells by $\|F_\rho\|_{L^1(P)}$, shows that
\[
 \frac1{|U_R|}\int\xi F_\rho\,\dd x
       \longrightarrow\frac1A\int_PF_\rho\,\dd x,
 \qquad
 \frac1{|U_R|}\sum_v\xi(v)\longrightarrow\frac NA.
\]
The resulting limit is
$(\int_PF_\rho+\pi N\log\rho)/A=\mathcal E_P(F)/A$.
\end{proof}

On $T_n=\R^2/(\sqrt n\,\Z^2)$, let $h$ be the mean-zero solution
of $-\Delta h=2\pi(\sum_{i=1}^n\delta_{v_i}-1)$ for distinct points
$v_1,\ldots,v_n$, with energy defined by \eqref{eq:torus-energy}.
The compatibility condition holds because $|T_n|=n$;
the mean-zero solution exists and is unique by the Fourier series
for the Laplacian on the torus.
Any periodic locally integrable curl-free field $E$ with divergence
$-\Delta h$ differs from
$-\nabla h$ by a smooth field of zero divergence and curl.
Each component is harmonic on the torus, hence constant.  Therefore
$E=-\nabla h+c$, with $c\in\R^2$.  As a distributional derivative
of a periodic function, $\nabla h$ has integral zero; it is in
$L^1(T_n)$ by Lemma~\ref{red:local}.  The cross integral on the
punctured torus thus tends to zero, giving the exact identity
\begin{equation}\label{red:harmonic}
 \mathcal E_{T_n}(E)=\mathcal W_n(h)+\frac n2|c|^2.
\end{equation}
Applying the same fixed orthogonal rotation to the fields preserves
this energy identity.

\begin{proposition}[Reduction to square tori]\label{red:reduction}
Let $C\in\R$.  Suppose that for every integer $n\ge2$ and every
configuration of $n$ distinct points of $T_n$, the mean-zero
potential in \eqref{eq:torus-energy} satisfies
$\mathcal W_n(h)\ge nC$.  Then $W(E)\ge C$ for every
$E\in\mathcal A_1$.  In particular, if these finite-torus bounds
hold with $C=W(E_\triangle)$, then
\[
 \inf_{E\in\mathcal A_1}W(E)=W(E_\triangle).
\]
\end{proposition}

\begin{proof}
Take the square-periodic minimizing sequence $j_k\in\mathcal C_1$
in \eqref{red:M}.  Integrating its curl on a square period torus
shows that the area of that torus is $2\pi N_k$, where $N_k$ is
its positive integer number of charges.  If $N_k=1$, use twice
the square period, which has four charges.  Lemma~\ref{red:periodic-average}
shows that this change of period preserves the energy per unit
area.  We may therefore assume $N_k\ge2$.  We may also align each
period square with the coordinate axes: simultaneous rotation of
the positions and vector components preserves its cell energy,
and the same lemma identifies its energy per unit area.

With $m=2\pi$, rescale and undo the component rotation:
\[
 E_k(x)=\sqrt m\bigl((j_k)_2(\sqrt m x),
                         -(j_k)_1(\sqrt m x)\bigr).
\]
This is an admissible field on $T_{N_k}$.  Changing the puncture
radius in its cell gives
\[
 \mathcal E_{T_{N_k}}(E_k)
 =\mathcal E_{P_k}(j_k)-\frac\pi2N_k\log m,
 \qquad |P_k|=mN_k.
\]
Combining periodic averaging, \eqref{red:harmonic}, and the
assumed torus inequality yields
\[
 m\left(\mathsf W_{\mathrm{ball}}(j_k)-\frac14\log m\right)
 =\frac{\mathcal E_{T_{N_k}}(E_k)}{N_k}\ge C.
\]
Letting $k\to\infty$ gives $m(M-\frac14\log m)\ge C$.
Lemma~\ref{red:plane-bound} now proves the conclusion for every
full-plane competitor, including all possible infinite-energy cases.

For the last assertion, the covolume-one triangular lattice is
simple, locally finite, and has quadratic charge growth.  Its
periodic potential gives an element $E_\triangle\in\mathcal A_1$,
whose finite energy follows from Lemmas~\ref{red:local} and
\ref{red:periodic-average}.  It supplies the opposite inequality
for the infimum.
\end{proof}
